% \paragraph{Context and motivation.}
Vehicular communication networks have undergone some important technological transitions. Initially, the focus was on Vehicular Ad hoc Networks (\acrotip{VANET}s), which prioritized direct connectivity between vehicles but faced challenges related to latency and connectivity due to the high mobility of nodes~\cite{kandali2025}. With the maturation of the Internet of Things (\acrotip{IoT}) and the expansion of embedded sensors in urban infrastructure, a large-scale network paradigm emerged that integrates vehicles, infrastructure, pedestrians, and the cloud, known as the Internet of Vehicles (\acrotip{IoV}) ecosystem~\cite{kandali2025}. This ecosystem was designed to enable truly autonomous and connected Intelligent Transportation Systems (\acrotip{ITS}), allowing vehicular applications to simultaneously meet efficiency, reliability, and safety requirements~\cite{uddin2025}.

However, \acrotip{IoV} depends on the ability to meet the low-latency and high-throughput requirements of radio technologies. Thus, vehicular communication has undergone several evolutions, starting initially with the allocation of \SI{75}{\mega\hertz} in the \SI{5{,}9}{\giga\hertz} band for the Dedicated Short-Range Communications (\acrotip{DSRC}) standard, which defined the parameters for Vehicle-to-Vehicle (\acrotip{V2V}) and Vehicle-to-Infrastructure (\acrotip{V2I}) communication in the early \acrotip{VANET}s~\cite{zheng2015}. From this legacy, two technological communication axes were developed. The first, \acrotip{IEEE}~802.11bd (Next-Generation \acrotip{V2X}), extended \acrotip{IEEE}~802.11p by incorporating robustness against fast fading and support for speeds up to \SI{250}{\kilo\meter\per\hour}~\cite{naik2019}. The second axis focuses on overcoming \acrotip{DSRC}'s limitations in scalability and deployment cost, which drove the development of Cellular \acrotip{V2X} (\acrotip{CV2X}) by the \acrotip{3GPP}~\cite{chen2017}. From a regulatory perspective, the dispute for the adoption of these technologies was resolved by the Federal Communications Commission (\acrotip{FCC}), which established \acrotip{CV2X} as the mandatory standard for \acrotip{ITS} communication, with the official discontinuation of \acrotip{DSRC} set for December 2026~\cite{fcc2024_59ghz}.

The evolution of communication technologies has succeeded in meeting the network latency and throughput requirements of \acrotip{ITS}. However, limitations regarding the computational capacity of embedded devices still needed to be overcome. In this context, Multi-Access Edge Computing (\acrotip{MEC}) and its vehicular specialization, Vehicular Edge Computing (\acrotip{VEC}), enable the operation of latency-sensitive applications by positioning processing and storage resources close to the radio access network~\cite{mao2017,miriyala2026dl_offload_review}. Thus, end users can transfer computationally intensive tasks to edge servers using \acrotip{5G} \acrotip{CV2X} infrastructure. This delegation drastically reduces response time and on-board energy consumption~\cite{liu2019,fang2021mec_iot,long2025}, thereby ensuring the Quality of Service (\acrotip{QoS}) of latency-sensitive applications.

In vehicular environments, the computational capacities of Interconnected On-Board Units (\acrotip{OBUI}s) are inherently heterogeneous~\cite{luo2024}. This hardware diversity, if well-utilized, enables the formation of cooperative execution environments, in which vehicles equipped with high-performance processing units make their temporary idle capacity available to assist neighboring nodes with computational and energy constraints. The direct sharing of these resources considerably improves the global \acrotip{QoS} of \acrotip{ITS}~\cite{he2019d2d_mec_twc}. Thus, integrating the high capacity of the edge server with the availability of idle processing from \acrotip{OBUI}s closes the circle of improvements necessary to enable the execution of computationally intensive, latency-sensitive, and energy-constrained vehicular applications.

The operationalization of this collaborative potential has task offloading as its central mechanism. This strategy consists of transferring the responsibility for computing complex tasks from resource-limited devices to more powerful remote processing nodes, such as edge servers or neighboring vehicles with greater processing capacity~\cite{wang2022}. The goal is to bypass local battery, storage, and processing constraints to reduce energy consumption, minimize latency, and improve \acrotip{QoS}. The great challenge, however, lies in dynamically deciding when and where to offload these tasks. Fully exploring this potential requires the development of allocation policies capable of adapting to the highly dynamic network topology and evaluating multiple decision criteria simultaneously and in real-time.

Task offloading in vehicular networks is fundamentally a multi-objective optimization problem~\cite{wang2020survey,chen2025}. In this context, various policies have been proposed to solve this problem. Heuristics based on logical rules and Multi-Criteria Decision Making (\acrotip{MCDM}) methods, such as greedy strategies, are widely used due to their computational simplicity and ease of implementation~\cite{ahmed2022survey}. However, these approaches only evaluate the instantaneous network situation, proving incapable of anticipating traffic dynamics and future states of the environment. On the other hand, meta-heuristics, such as genetic algorithms, particle swarm optimization, and ant colony optimization, offer more solid global search guarantees and have been successfully applied to offloading decision-making in static scenarios~\cite{li2020genetic,you2021efficient,kishor2022task}. However, these methods may prove inadequate for vehicular environments, as population-based search requires tens to hundreds of iterative fitness evaluations at each decision point, which can cause unacceptable delays~\cite{wu2024,cao2024survey}. This relatively slow convergence characteristic can be incompatible with the rigorous execution deadlines on the millisecond scale characteristic of vehicular applications~\cite{miriyala2026dl_offload_review,uddin2025}. Furthermore, these algorithms need to be restarted with each topological change in the environment, and their initialization and convergence complexity grows significantly with the increase in network size, potentially degrading the quality of the solution in large-scale scenarios~\cite{wang2020survey,long2025}.

Simplified evaluation scenarios, with a fixed number of vehicles, uniform task distributions, and idealized channel models, do not capture the complexity and high dynamism of real vehicular networks~\cite{miriyala2026dl_offload_review}. Offloading policies trained and evaluated in these environments can suffer performance degradation when vehicular density or workload distribution changes suddenly~\cite{miriyala2026dl_offload_review}, revealing a critical dichotomy between the success obtained in simulations and real applicability in operational scenarios. In this context, it is fundamental that the development of robust offloading policies actively seeks to mitigate this disparity. This requires the design of adaptive state representations that ensure the model's generalization capability in the face of constant changes in the number of candidate nodes for offloading, eliminating the dependence on stationary traffic environments~\cite{shao2025,uddin2025}.

\acrotip{DRL} presents itself as a natural solution for task offloading because it learns a control policy that maps observations to actions considering long-term future consequences, rather than just the instantaneous state of the network~\cite{liu2019,shi2023}. Unlike heuristics and meta-heuristics, a \acrotip{DRL} agent has the ability to explore temporal correlations in traffic patterns, queue dynamics, vehicular mobility, and communication channel quality from continuous interaction with the environment~\cite{li2025drl_survey,uddin2025}.

However, standard \acrotip{DRL} architectures share a critical structural weakness related to their practical applicability, as they frequently require retraining whenever the operational scenario undergoes changes in the number of candidate computational nodes~\cite{miriyala2026dl_offload_review}. Since the input layer of deep neural networks depends on a fixed-size vector, constant changes in the number of candidate nodes for offloading, such as the edge server and neighboring vehicles, end up invalidating the previously trained model~\cite{uddin2025}.

In this context, this work proposes a hybrid task offloading policy called \acrotip{SAC}-\acrotip{TOPSIS} capable of handling the operational constraints of \acrotip{DRL} in a dynamic vehicular environment. In summary, the main contributions are:

\begin{enumerate}[leftmargin=*,nosep]
  \item Development of an innovative solution for task offloading that combines the \acrotip{TOPSIS} technique with the \acrotip{SAC} algorithm. \acrotip{TOPSIS} provides the best alternative based on instantaneous network conditions, while \acrotip{SAC} learns the long-term optimal policy through accumulated experience.
  \item Representation of a variable-sized set of candidate nodes for offloading in a compact six-dimensional state vector, allowing the control policy trained in a representative scenario to be tested in unseen environments without the need for retraining.
  \item Reduction of the action space to a binary decision (local processing vs. best offloading candidate), enabling inference in $O(1)$ direct passes through the neural network per task, making the model suitable for real-time embedded deployment.
  \item Development of a deadline-sensitive reward function that provides a continuous gradient signal in both success and failure regimes, avoiding reward sparsity and teaching the agent to actively minimize delay rather than merely avoid it.
\end{enumerate}

The remainder of this article is organized as follows. Section~\ref{sec:related} reviews related work in the literature. Section~\ref{sec:model} formalizes the system model and defines the optimization objective. Section~\ref{sec:arch} details the proposed \acrotip{SAC}-\acrotip{TOPSIS} hybrid policy architecture. Section~\ref{sec:eval} presents the experimental evaluation and discusses the results obtained. Section~\ref{sec:conclusion} concludes the article, presenting final considerations and directions for future work. Furthermore, the definitions of all acronyms used throughout the text are summarized in Table~\ref{tab:acronyms}.

\begin{table*}[t]
\caption{List of Acronyms -- Summary of the main acronyms used throughout this article.}
\label{tab:acronyms}
{\scriptsize
\begin{tabularx}{\textwidth}{l X l X}
\toprule
\textbf{Acronym} & \textbf{Definition} & \textbf{Acronym} & \textbf{Definition} \\
\midrule
\acrotip{3GPP}   & 3rd Generation Partnership Project & \acrotip{NIS}    & Negative Ideal Solution \\
\acrotip{5G}     & 5th Generation                      & \acrotip{NRV2X}  & New Radio Vehicle-to-Everything \\
\acrotip{CV2X}   & Cellular Vehicle-to-Everything      & \acrotip{OBUI}   & Interconnected On-Board Unit \\
\acrotip{D3QN}   & Dueling Double Deep Q-Network       & \acrotip{PIS}    & Positive Ideal Solution \\
\acrotip{DRL}    & Deep Reinforcement Learning         & \acrotip{QoS}    & Quality of Service \\
\acrotip{DSRC}   & Dedicated Short-Range Comm.         & \acrotip{RSU}    & Roadside Unit \\
\acrotip{IEEE}   & Inst. of Elect. and Electr. Eng.    & \acrotip{SAC}    & Soft Actor-Critic \\
\acrotip{IoT}    & Internet of Things                  & \acrotip{TOPSIS} & Tech. Order Pref. Sim. Ideal Sol. \\
\acrotip{IoV}    & Internet of Vehicles                & \acrotip{V2I}    & Vehicle-to-Infrastructure \\
\acrotip{ITS}    & Intelligent Transportation Systems  & \acrotip{V2N}    & Vehicle-to-Network \\
\acrotip{JCT}    & Job Completion Time                 & \acrotip{V2P}    & Vehicle-to-Pedestrian \\
\acrotip{MCDM}   & Multi-Criteria Decision Making      & \acrotip{V2V}    & Vehicle-to-Vehicle \\
\acrotip{MEC}    & Multi-Access Edge Computing         & \acrotip{VEC}    & Vehicular Edge Computing \\
\bottomrule
\end{tabularx}}
\end{table*}
